\documentclass[a4paper,11pt]{article}
\usepackage[utf8]{inputenc}
\usepackage[T1]{fontenc}
\usepackage[ngerman]{babel}
\usepackage[margin=2cm]{geometry}
\usepackage{qrcode}
\usepackage{hyperref}
\usepackage{xcolor}

\definecolor{physik}{HTML}{2c5aa0}

\pagestyle{empty}

\begin{document}

\begin{center}
\vspace*{2cm}

{\Huge\bfseries\color{physik} Lernpfad}

\vspace{0.5cm}

{\Large Bohr-Atommodell}

\vspace{0.3cm}

{\normalsize Physik 12 GK | Mi 28.01. + Fr 30.01.2026}

\vspace{2cm}

\qrcode[height=6cm]{https://devbox42.github.io/sciencesim/physik/kl12-GK/01-photoeffekt/06-ES-bohr-postulate/LP-06-bohr-postulate.html}

\vspace{1.5cm}

{\small\color{gray}
\url{https://devbox42.github.io/sciencesim/physik/kl12-GK/01-photoeffekt/06-ES-bohr-postulate/LP-06-bohr-postulate.html}
}

\vspace{2cm}

\begin{tabular}{|l|l|}
\hline
\textbf{Thema} & Bohr-Postulate, Energieniveaus, Balmer-Serie \\
\hline
\textbf{Dauer} & Mi 45 min + Fr 90 min \\
\hline
\textbf{Material} & Arbeitsblatt AB-06 (15 BE) \\
\hline
\end{tabular}

\end{center}

\end{document}
